\Askhsh{16097}{1}
\begin{ylh}{1}
    \item 7.7. Θεώρημα του Θαλή
    \item 8.2. Κριτήρια ομοιότητας
    \item 9.1. Ορθές προβολές
    \item 10.3. Εμβαδόν βασικών ευθύγραμμων σχημάτων
    \item 11.2. Ιδιότητες και στοιχεία κανονικών πολυγώνων.
\end{ylh}
\begin{wrapfigure}[7]{r}{4.5cm}
    \begin{tikzpicture}
        % Define points of the right triangle
        \tkzDefPoint(0,0){A}
        \tkzDefPoint(4,0){G} % Gamma
        \tkzDefPoint(0,3){B}

        % Draw the triangle
        \tkzDrawPolygon(A,B,G)

        % Draw the right angle at A
        \tkzMarkRightAngle(G,A,B)

        % Define the foot of the altitude D from A to BG
        \tkzDefPointBy[projection=onto line G--B](A) \tkzGetPoint{D}

        % Draw the altitude AD
        \tkzDrawSegment(A,D)

        % Draw the right angle at D
        \tkzMarkRightAngle(A,D,G)

        % Label the points
        \tkzLabelPoint[below left](A){A}
        \tkzLabelPoint[above left](B){B}
        \tkzLabelPoint[below right](G){$\Gamma$}
        \tkzLabelPoint[above right](D){$\Delta$}
    \end{tikzpicture}
\end{wrapfigure}
\begin{alist}
    \item Να χαρακτηρίσετε καθεμία από τις προτάσεις που ακολουθούν ως \textbf{Σωστή (Σ)} ή \textbf{Λανθασμένη (Λ)}, γράφοντας στην κόλλα σας, δίπλα στον αριθμό που αντιστοιχεί σε καθεμιά από αυτές το γράμμα Σ αν η πρόταση είναι Σωστή, ή το γράμμα Λ αν αυτή είναι Λάθος.
    \begin{rlist}
        \item Κάθε ευθεία που είναι παράλληλη με μία από τις πλευρές ενός τριγώνου χωρίζει τις δύο άλλες πλευρές σε μέρη ανάλογα.
        \item Δύο ισοσκελή τρίγωνα είναι πάντοτε όμοια.
        \item Στο σχήμα, η προβολή της πλευράς $AB$ στην υποτείνουσα $B\Gamma$ είναι το τμήμα $\Gamma\Delta$.
        \item Το εμβαδόν ενός τριγώνου ισούται με το γινόμενο μιας πλευράς επί το αντίστοιχο ύψος.
        \item Ο εγγεγραμμένος και ο περιγεγραμμένος κύκλος ενός κανονικού πολυγώνου είναι ομόκεντροι.
    \end{rlist} \monades{10}
    \item Να αποδείξετε ότι το εμβαδόν ενός ορθογωνίου παραλληλογράμμου ισούται με το γινόμενο των πλευρών του. \monades{15}
\end{alist}